\documentclass[11pt,a4paper]{article}

\usepackage[utf8]{inputenc}
\usepackage[T1]{fontenc}
\usepackage[spanish,es-noquoting]{babel}
\usepackage{amsmath,amssymb,amsthm}
\usepackage{booktabs}
\usepackage{siunitx}
\usepackage[margin=2.6cm]{geometry}
\usepackage{xcolor}
\usepackage{hyperref}
\hypersetup{colorlinks=true,linkcolor=blue!50!black,citecolor=blue!50!black,urlcolor=blue!50!black}

\newcommand{\ux}{\hat{\mathbf{u}}}
\newcommand{\nhat}{\hat{\mathbf{N}}}
\newcommand{\rhat}{\hat{\mathbf{r}}}
\newcommand{\phihat}{\hat{\boldsymbol{\varphi}}}
\newcommand{\zhat}{\hat{\mathbf{z}}}
\newcommand{\thhat}{\hat{\boldsymbol{\theta}}}
\newcommand{\E}{\mathbf{E}}

\title{Refracción en un sistema estigmático simple:\\
transferencia vectorial, flujo energético y reducción al plano focal}
\author{Jhonatan S. Blanco}
\date{Versión verificada numéricamente}

\begin{document}
\maketitle

\begin{abstract}
Se deriva la transferencia vectorial y energética de un campo electromagnético
monocromático a través de un dioptrio estigmático simple que separa dos medios
dieléctricos, homogéneos, isótropos y no magnéticos, de índices $n_0$ y $n_i$.
El campo incidente se prescribe sobre una esfera $\mathcal{G}$ centrada en el
punto objeto $A$; el transmitido se evalúa sobre una esfera $\mathcal{G}'$
centrada en el punto imagen $A'$. La conservación del flujo medio de Poynting
determina el factor geométrico de amplitud
\begin{equation}
  \mathcal{A}(Q) \;=\; \frac{|z_0|}{|z_i|}\,\frac{\ell_i(Q)}{\ell_0(Q)} .
  \label{eq:A}
\end{equation}
Respecto de versiones anteriores de esta nota, aquí se corrigen o se agregan
cinco puntos: \emph{(i)} la superficie se parametriza por el radio cilíndrico
$r$ del punto $Q$, que es la coordenada de apertura, y no por el radio
esférico $\rho$ desde el vértice; \emph{(ii)} se caracterizan los dos
mecanismos distintos que acotan la apertura útil ---incidencia rasante y
ángulo crítico---; \emph{(iii)} el paso de $\mathcal{G}'$ al plano focal se
identifica como una reducción de Debye, cuyo cambio de variable es el
\emph{mapeo seno} $u=|z_i|\sin\alpha_i$ con jacobiano $1/\cos\alpha_i$, y no
el mapeo perspectivo; \emph{(iv)} se establece cuál es la versión escalar
correcta del problema, que resulta ser el canal $s$; \emph{(v)} todos los
resultados se contrastan contra dos referencias exactas independientes ---la
integral de Franz para Maxwell y la de Helmholtz--Kirchhoff para el problema
escalar---, con acuerdo a $10^{-8}$ en amplitud absoluta y fase.
\end{abstract}

\section{Geometría del sistema}

Sea $\Sigma\subset\mathbb{R}^3$ una superficie refractora de revolución
alrededor del eje $z$, con vértice en el origen $O$. El sistema es
exactamente estigmático para el par
\begin{equation}
  A=(0,0,z_0),\qquad A'=(0,0,z_i),\qquad z_0,z_i\in\mathbb{R}\setminus\{0\},
\end{equation}
y las esferas de referencia tangentes al plano $z=0$ en $O$ son
\begin{equation}
  \mathcal{G}=\{X: \|X-A\|=|z_0|\},\qquad
  \mathcal{G}'=\{X: \|X-A'\|=|z_i|\}.
\end{equation}

\subsection{Parametrización: el radio cilíndrico, no el esférico}
\label{sec:param}

Un punto $Q\in\Sigma$ admite dos radios naturales: el cilíndrico
$r=\sqrt{Q_x^2+Q_y^2}$ y el esférico desde el vértice
$\rho=\|Q-O\|$, ligados por
\begin{equation}
  \rho^2 = r^2 + z(r)^2 .
  \label{eq:rho_r}
\end{equation}
La forma cerrada del óvalo cartesiano se escribe naturalmente en $\rho$, pero
\textbf{la coordenada física de apertura es $r$}: es el radio del diafragma
sobre el dioptrio y el que fija sin ambigüedad el rayo. Confundirlas no es un
detalle de notación. Para $n_0=1$, $n_i=1{,}5$, $z_0=\SI{-10}{\milli\meter}$,
$z_i=\SI{6}{\milli\meter}$, la incidencia rasante ocurre en $r=2{,}397$ pero
en $\rho=3{,}758$: una apertura citada como fracción de $\rho$ se sale de la
superficie útil en un $57\%$.

La inversión $r\mapsto\rho$ requiere resolver una cuártica. En la práctica
conviene hacerla numéricamente ---tabla monótona densa refinada por Newton
sobre $r(\rho)-r$---, lo que alcanza precisión de máquina ($2\times10^{-15}$
en las verificaciones de la Sec.~\ref{sec:verif}).

\subsection{Distancias y ángulos}

Para cada $Q$ se definen las distancias positivas
$\ell_0(Q)=\|Q-A\|$, $\ell_i(Q)=\|A'-Q\|$, y con
$z\equiv z(r)$,
\begin{equation}
  \ell_0=\sqrt{r^2+[z-z_0]^2},\qquad
  \ell_i=\sqrt{r^2+[z_i-z]^2}.
\end{equation}
Las direcciones físicas de propagación son
$\ux_0=(Q-A)/\ell_0$ y $\ux_i=(A'-Q)/\ell_i$, que en la base cilíndrica
$(\rhat,\phihat,\zhat)$ se escriben
\begin{equation}
  \ux_0=\sin\alpha_0\,\rhat+\cos\alpha_0\,\zhat,\qquad
  \ux_i=-\sin\alpha_i\,\rhat+\cos\alpha_i\,\zhat,
\end{equation}
con
\begin{equation}
  \sin\alpha_0=\frac{r}{\ell_0},\quad
  \cos\alpha_0=\frac{z-z_0}{\ell_0},\qquad
  \sin\alpha_i=\frac{r}{\ell_i},\quad
  \cos\alpha_i=\frac{z_i-z}{\ell_i}.
\end{equation}
Los senos son no negativos para $r\ge0$; los cosenos conservan el signo de la
proyección axial, lo que describe las distintas posiciones de $A$ y $A'$ sin
fijar de antemano los signos de $z_0$ y $z_i$.

Sea $\nhat$ la normal unitaria a $\Sigma$ orientada hacia el medio incidente.
Los cosenos de incidencia y transmisión son
\begin{equation}
  \cos\theta_0=-\ux_0\cdot\nhat,\qquad \cos\theta_i=-\ux_i\cdot\nhat,
  \label{eq:cosines}
\end{equation}
y rige la ley de Snell $n_0\sin\theta_0=n_i\sin\theta_i$.

\begin{quote}\small
\textbf{Nota de implementación.} La orientación de $\nhat$ debe fijarse de
forma global (por el lado en que está $A$), no por un test construido con el
rayo incidente: cualquier criterio del tipo $\operatorname{sign}(\ux_0\cdot\nhat)$
degenera exactamente en la rasancia, que es el único lugar donde importa.
\end{quote}

La condición estigmática, dado que $O\in\Sigma$ y las esferas pasan por $O$,
adopta la forma
\begin{equation}
  n_0\,\ell_0(Q)+n_i\,\ell_i(Q)=n_0|z_0|+n_i|z_i|,\qquad Q\in\Sigma .
  \label{eq:stigmatic}
\end{equation}

\section{Límites de la apertura útil}
\label{sec:apertura}

La superficie no se extiende indefinidamente: hay dos mecanismos distintos que
la acotan, y cuál manda depende del contraste de índices.

\paragraph{Rasancia ($n_0<n_i$).} El ángulo de incidencia crece con $r$ y
alcanza $90^\circ$ a radio finito: $\cos\theta_0\to0$. Más allá, la forma
cerrada continúa sobre una hoja del óvalo que $A$ ya no ilumina.

\paragraph{Ángulo crítico ($n_0>n_i$).} Aquí es el ángulo \emph{transmitido}
el que alcanza $90^\circ$ primero, en $\sin\theta_0=n_i/n_0$. Pasado ese
radio la ley de Snell no tiene solución real: \textbf{no existe óvalo} que
envíe el rayo al punto imagen. El límite es intrínseco y no se corre eligiendo
conjugados.

Esta asimetría tiene una consecuencia de diseño fuerte. Una superficie de
salida denso$\to$raro satura: con sílice a aire, la apertura numérica imagen no
pasa de $\mathrm{NA}\approx0{,}74$ aunque el objeto intermedio se aleje
arbitrariamente. Superarlo exige subir el índice del último elemento o
introducir un medio de inmersión ---que es la ruta que tomó la litografía de
alta apertura.

Además, cerca del ángulo crítico los coeficientes de Fresnel permanecen
finitos ($t_s,t_p\to$ valores del orden de $2$) mientras
$\mathcal{A}$ crece sin freno, de modo que el peso de pupila
\emph{diverge} hacia el borde. Para superficies raro$\to$denso ocurre lo
contrario: $t_{s,p}\to0$ en la rasancia y la pupila se apaga suavemente. Es
una diferencia cualitativa entre ambos tipos de superficie, no un matiz.

\section{Bases esféricas y flujo en un tubo de rayos}

La base esférica incidente centrada en $A$ es $(\ux_0,\thhat_0,\phihat)$ con
$\thhat_0=\phihat\times\ux_0=\cos\alpha_0\,\rhat-\sin\alpha_0\,\zhat$, y la
final centrada en $A'$ es $(\ux_i,\thhat_i,\phihat)$ con
$\thhat_i=\cos\alpha_i\,\rhat+\sin\alpha_i\,\zhat$. En una superficie de
revolución $\phihat$ coincide con la dirección $s$ de Fresnel y $\thhat_{0,i}$
son las direcciones meridionales $p$. La transversalidad deja dos amplitudes
independientes sobre cada esfera.

Con dependencia temporal $e^{-i\omega t}$ y $k=2\pi/\lambda$ en el vacío, una
onda localmente plana y transversal en un medio de índice real $n$ transporta
$\langle\mathbf{S}\rangle=(n/2Z_{\rm vac})\|\E\|^2\ux$. Igualando la potencia
sobre $\mathcal{G}$, sobre la interfaz y sobre $\mathcal{G}'$, y usando las
identidades geométricas $\cos\theta_0\,dS_\Sigma=\ell_0^2\,d\Omega_0$ y
$\cos\theta_i\,dS_\Sigma=\ell_i^2\,d\Omega_i$, se obtiene la transferencia
completa
\begin{equation}
  \E_{\mathcal{G}'}
  =\mathcal{A}(Q)\Bigl[t_p(Q)\,E_{\mathcal{G}}^{\theta_0}\,\thhat_i
  +t_s(Q)\,E_{\mathcal{G}}^{\varphi}\,\phihat\Bigr],
  \label{eq:transfer}
\end{equation}
con $\mathcal{A}$ dada por \eqref{eq:A} y los coeficientes dieléctricos
\begin{equation}
  t_s=\frac{2n_0\cos\theta_0}{n_0\cos\theta_0+n_i\cos\theta_i},\qquad
  t_p=\frac{2n_0\cos\theta_0}{n_i\cos\theta_0+n_0\cos\theta_i}.
  \label{eq:fresnel}
\end{equation}
La condición estigmática elimina la fase geométrica relativa entre las dos
esferas. Las transmitancias de potencia asociadas,
\begin{equation}
  T_q=\frac{n_i\cos\theta_i}{n_0\cos\theta_0}|t_q|^2,\qquad R_q+T_q=1,
  \label{eq:T}
\end{equation}
verifican el balance
\begin{equation}
  \frac{d\Phi_{\mathcal{G}',q}}{d\Phi_{\mathcal{G},q}}=T_q ,
  \label{eq:balance}
\end{equation}
que es la identidad que fija $\mathcal{A}$ de forma independiente de todo lo
demás.

\section{Matriz de transferencia}

Eliminando la componente longitudinal por transversalidad sobre $\mathcal{G}$
($E_{\mathcal{G}}^{\theta_0}=E_{\mathcal{G}}^{r}/\cos\alpha_0$) se obtiene la
forma reducida, diagonal en la base cilíndrica:
\begin{equation}
  \lambda_r=\mathcal{A}\,t_p\,\frac{\cos\alpha_i}{\cos\alpha_0},\qquad
  \lambda_\varphi=\mathcal{A}\,t_s,\qquad
  \lambda_z=\mathcal{A}\,t_p\,\frac{\sin\alpha_i}{\cos\alpha_0}.
  \label{eq:eigen}
\end{equation}
En la base cartesiana, con $\phi$ el azimut de pupila,
\begin{equation}
  \mathsf{T}_\perp=
  \frac{\lambda_r+\lambda_\varphi}{2}\,\sigma_0
  +\frac{\lambda_r-\lambda_\varphi}{2}
  \bigl(\cos2\phi\,\sigma_1+\sin2\phi\,\sigma_2\bigr),
  \label{eq:cartesian}
\end{equation}
y $E_{\mathcal{G}'}^{z}=\lambda_z\,(E_{\mathcal{G}}^{x}\cos\phi+E_{\mathcal{G}}^{y}\sin\phi)$.

\paragraph{Generalidad.} Toda la derivación es \emph{puntual}: tubo de rayos
local y reparto de Fresnel local en $Q$. En ningún paso interviene la
estructura azimutal de $\E_{\mathcal{G}}$, de modo que \eqref{eq:transfer} vale
para campo incidente arbitrario ---cualquier polarización, cualquier
dependencia en $\varphi$. Lo único atado a la superficie de revolución es la
identificación $s\equiv\phihat$, $p\equiv\thhat$. Para una superficie general
basta reemplazar el proveedor de marcos locales; la física no cambia.

\paragraph{Transversalidad.} De \eqref{eq:eigen} se sigue
\begin{equation}
  \frac{\lambda_z}{\lambda_r}=\tan\alpha_i ,
  \label{eq:transversality}
\end{equation}
que es exactamente la condición $\mathbf{k}\cdot\E=0$ para cada modo del
espectro angular. La transversalidad sobre $\mathcal{G}'$ y la condición de
divergencia nula en el espacio espectral son el mismo enunciado.

\section{De $\mathcal{G}'$ al plano focal: la reducción de Debye}
\label{sec:debye}

Esta sección reemplaza el tratamiento perspectivo de versiones anteriores.

El campo en el plano focal es la integral de Debye sobre la esfera de salida,
\begin{equation}
  \E(P)=e^{ik_i|z_i|}\,\frac{-ik_i|z_i|}{2\pi}
  \int\!\!\int \E_{\mathcal{G}'}(\hat{\mathbf{s}})\,
  e^{ik_i\hat{\mathbf{s}}\cdot\mathbf{r}_P}\,d\Omega ,
  \qquad k_i=\frac{2\pi n_i}{\lambda}.
  \label{eq:debye}
\end{equation}
Es una superposición de ondas planas con $|\mathbf{k}|=n_ik_0$ \emph{exacto}:
espectro angular, no una aproximación de Fresnel.

Sustituyendo el \textbf{mapeo seno}
\begin{equation}
  u=|z_i|\sin\alpha_i,\qquad
  d\Omega=\frac{u\,du\,d\phi}{|z_i|^2\cos\alpha_i},
  \label{eq:sine}
\end{equation}
la integral queda
\begin{equation}
  \E(P)=e^{ik_i|z_i|}\,\frac{-ik_i}{2\pi|z_i|}
  \int\!\!\int \frac{\E_{\mathcal{G}'}}{\cos\alpha_i}\,
  e^{iq\,u\cos(\phi-\phi_P)}\,u\,du\,d\phi ,
  \qquad q=\frac{k_i\rho_P}{|z_i|}.
\end{equation}
Es decir: \textbf{la transformada integra sobre $u=|z_i|\sin\alpha_i$ y el
integrando lleva el jacobiano $1/\cos\alpha_i$}. La coordenada focal resulta
$\rho_P=\lambda z_i q/(2\pi n_i)$, y el prefactor absoluto es
\begin{equation}
  \mathcal{P}=e^{ik_i|z_i|}\,\frac{-ik_i}{|z_i|} ,
  \label{eq:prefactor}
\end{equation}
una vez absorbido el $2\pi$ de la reducción azimutal.

\paragraph{Qué estaba mal.} El mapeo perspectivo $r_P=|z_i|\tan\alpha_i$ con
peso $\cos\alpha_i$ es \emph{correcto como relación de flujo entre el plano
tangente y la esfera} ---el campo sobre el plano realmente es $\cos\alpha_i$
veces el campo sobre la esfera---, pero \textbf{no es el cambio de variable de
esta transformada}. Usarlo cuesta un $14\%$ de error en amplitud focal a
apertura alta. La distinción es, en el fondo, la misma que separa espectro
angular de Fresnel: en el borde de una pupila con
$\sin\alpha_i=0{,}614$, los dos mapeos difieren un $26{,}7\%$
($\tan$ contra $\sin$).

\section{El problema escalar}
\label{sec:escalar}

¿Cuál es la versión escalar correcta de este dioptrio? El problema escalar bien
planteado es la ecuación de Helmholtz en ambos medios con $U$ y $\partial U/\partial n$
continuas sobre $\Sigma$. Para una onda localmente plana, esa continuidad da
como coeficiente de transmisión
\begin{equation}
  t=\frac{2n_0\cos\theta_0}{n_0\cos\theta_0+n_i\cos\theta_i}=t_s ,
\end{equation}
es decir, \textbf{el problema escalar es el problema TE}. Y como el transporte
de flujo en el tubo de rayos es ciego a la polarización, $\mathcal{A}(Q)$ y el
mapeo seno se trasladan sin cambio. Por lo tanto:

\begin{equation}
  \boxed{\;\text{versión escalar correcta}\;=\;
  \text{canal } s:\quad w_s=\mathcal{A}\,t_s\,\frac{1}{\cos\alpha_i}\;}
\end{equation}

sin mezcla de polarización y sin $E_z$. Conviene subrayar dos alternativas que
\emph{no} lo son:
\begin{itemize}
  \item el promedio $(\lambda_r+\lambda_\varphi)/2$ arrastra la proyección
    meridional $\cos\alpha_i/\cos\alpha_0$, que es un objeto vectorial sin
    lugar en un problema escalar;
  \item la pupila plana $t\equiv1$ no resuelve ninguna interfaz: es un modelo
    de lente ideal, útil como referencia de resolución pero no como solución.
\end{itemize}

\section{Verificación numérica}
\label{sec:verif}

Todos los resultados anteriores se contrastaron contra dos referencias
exactas construidas de forma independiente del modelo de rayos.

\paragraph{Árbitro vectorial.} La integral de Franz/Stratton--Chu sobre la
superficie real, con datos de contorno de óptica física. Es un campo de
Maxwell exacto en el semiespacio imagen. Calibración: reproduce el patrón de
Airy a $1{,}1\times10^{-3}$ y la amplitud de Debye con residuo
$\to \mathrm{NA}^2/8$ exacto al bajar la apertura ($0{,}1295\to0{,}1251$).

\paragraph{Árbitro escalar.} La integral de Helmholtz--Kirchhoff sobre la misma
superficie, con $U^+=t_s\,U^-$ y
$\partial U/\partial N=-ik_i\cos\theta_i\,U^+$. Calibración: $1{,}000000$
contra el caso analítico de onda convergente.

\begin{table}[htbp]
\centering
\small
\begin{tabular}{lll}
\toprule
Cantidad verificada & Criterio & Residuo \\
\midrule
Inversión $r\leftrightarrow\rho$ & $r(\rho(r))=r$ & $2\times10^{-15}$ \\
Condición estigmática \eqref{eq:stigmatic} & constante sobre $\Sigma$ & $4\times10^{-15}$ \\
Ley de Snell & $n_0\sin\theta_0=n_i\sin\theta_i$ & $3\times10^{-13}$ \\
Conservación de Fresnel & $R_q+T_q=1$ & $5\times10^{-16}$ \\
Balance de flujo \eqref{eq:balance} & $d\Phi'/d\Phi=T_q$ & $10^{-10}$ \\
Transversalidad \eqref{eq:transversality} & $\mathbf{k}\cdot\E=0$ & $4\times10^{-11}$ \\
Ecuación de onda & residuo de Helmholtz & $O(\delta z^2)$ puro \\
\midrule
\textbf{Campo focal vectorial} & vs.\ Franz, amplitud absoluta & $1{,}000000$ \\
 & vs.\ Franz, fase & $+0{,}0002^\circ$ \\
 & vs.\ Franz, RMS del perfil & $1{,}4\times10^{-8}$ \\
 & canal longitudinal (orden 1) & $1{,}000000$ \\
\midrule
\textbf{Campo focal escalar} & canal $s$ vs.\ Kirchhoff & $1{,}000000$ \\
 & promedio $(\lambda_r+\lambda_\varphi)/2$ & $1{,}0246$ \\
 & pupila plana $t\equiv1$ & $0{,}3040$ \\
\midrule
\textbf{Mapeo de pupila} & seno vs.\ Franz (RMS) & $2\times10^{-8}$ \\
 & tangente vs.\ Franz (RMS) & $4\times10^{-2}$ \\
\bottomrule
\end{tabular}
\caption{Verificación de la nota. Las tres primeras filas son identidades
geométricas; las siguientes, físicas; el bloque final contrasta el modelo
contra las dos referencias exactas.}
\end{table}

Vale subrayar el renglón de la ecuación de onda: el residuo de
$(\partial_z^2+\nabla_\perp^2+n^2k_0^2)U$ cae en razón exactamente $4{,}00$ al
halvar el paso del stencil, o sea que \emph{todo} el residuo es error de
discretización de la medición y ninguno del campo.

\section{Alcance y límites}

Lo que este desarrollo \textbf{sí} garantiza: dado el campo sobre
$\mathcal{G}$, el campo en el plano focal es solución exacta de Maxwell (caso
vectorial) o de Helmholtz (caso escalar), con amplitud y fase absolutas
correctas, sin aproximación paraxial ni de Fresnel en ningún punto de la
cadena.

Lo que \textbf{no} elimina:
\begin{itemize}
  \item \emph{La condición de contorno.} El reparto de Fresnel local en $Q$ es
    óptica física. Las dos referencias comparten esa hipótesis, así que la
    verificación de la Sec.~\ref{sec:verif} no la pone a prueba. Levantarla
    exige un solve de contorno riguroso (Müller/PMCHWT sobre cuerpo de
    revolución).
  \item \emph{Estigmatismo puntual.} El óvalo es perfecto para \emph{un} par
    de puntos conjugados, el axial. Fuera de eje aberra, y a apertura alta
    rápido: se midió que la PSF deja de estar limitada por difracción más allá
    de unas cinco longitudes de onda de altura de objeto. Un óvalo aislado es
    un formador de puntos, no un sistema de campo.
  \item \emph{Máscaras delgadas.} Si el objeto es una transmitancia
    $T(x,y)$, se supone $\E=T\,\E_{\rm inc}$ en su plano (Kirchhoff de máscara
    delgada), razonable mientras las estructuras sean grandes frente a
    $\lambda$.
\end{itemize}

\section{Conclusión}

La refracción en un sistema estigmático simple queda descrita, en las bases
esféricas centradas en los focos, por \eqref{eq:transfer}. El factor
$\mathcal{A}=|z_0|\ell_i/(|z_i|\ell_0)$ representa la redistribución
geométrica de amplitud entre las dos esferas y vale uno en el eje: es una
apodización de apertura alta pura. Los coeficientes $t_s$ y $t_p$ fijan la
amplitud de la rama transmitida y sus transmitancias satisfacen
\eqref{eq:T}--\eqref{eq:balance}.

El paso al plano focal es una reducción de Debye sobre el mapeo seno, con el
prefactor absoluto \eqref{eq:prefactor}; el mapeo perspectivo, aunque correcto
como relación de flujo plano--esfera, no es el cambio de variable de esa
transformada. La versión escalar del problema es el canal $s$, porque la
continuidad de $U$ y $\partial U/\partial n$ reproduce exactamente $t_s$.

Finalmente, la identidad $\lambda_z=\lambda_r\tan\alpha_i$ no es un peso
auxiliar sino la condición de divergencia nula expresada rayo a rayo, lo que
explica que el campo construido sea exactamente transversal a precisión de
máquina.

\begin{thebibliography}{9}
\bibitem{bornwolf} M. Born y E. Wolf, \emph{Principles of Optics}, 7.\textsuperscript{a} ed., Cambridge University Press, 1999.
\bibitem{jackson} J. D. Jackson, \emph{Classical Electrodynamics}, 3.\textsuperscript{a} ed., Wiley, 1999.
\bibitem{richardswolf} B. Richards y E. Wolf, ``Electromagnetic diffraction in optical systems. II. Structure of the image field in an aplanatic system'', \emph{Proc.\ R.\ Soc.\ A} \textbf{253}, 358--379 (1959).
\bibitem{goodman} J. W. Goodman, \emph{Introduction to Fourier Optics}, 4.\textsuperscript{a} ed., W. H. Freeman, 2017.
\bibitem{stratton} J. A. Stratton y L. J. Chu, ``Diffraction theory of electromagnetic waves'', \emph{Phys.\ Rev.} \textbf{56}, 99--107 (1939).
\bibitem{flagello} D. G. Flagello, T. Milster y A. E. Rosenbluth, ``Theory of high-NA imaging in homogeneous thin films'', \emph{J.\ Opt.\ Soc.\ Am.\ A} \textbf{13}, 53--64 (1996).
\end{thebibliography}

\end{document}
